\section{Conclusions（结论）}
\label{sec:conclusions}

本文提出神经增广 POD–Galerkin 方法（NE-PODG），用于参数化稳态传热问题的投影降阶。以训练快照的场值监督训练的网络，在查询参数处给出一组候选场；候选场去除 SA-PODG 基础空间的分量并正交化后构成神经互补空间，与基础空间一起组成增广试探空间，降阶解由该空间上的 Galerkin 闭合确定。

在六个参数化稳态传热问题上，NE-PODG 降低了 SA-PODG 的主场误差与传热量误差，问题覆盖线性对流扩散、共轭传热、自然对流、混合对流、多孔介质与开口分离流动。逐工况误差比的中位数为 2.2--27，其中五个算例只用 4 个快照。增益在快照稀疏时最大：方腔自然对流中，NE-PODG 用 12 个快照的外推精度与 SA-PODG 用 30 个快照相当。

精度收益主要来自神经互补空间对基础空间余量的捕获。训练网络相对同结构的未训练网络，把 $1-\Gamma_\theta$ 降低一个数量级以上；跨算例的增益次序与 $\Gamma_\theta$ 一致，$\chi_B/\chi_W$ 在方腔类算例上接近 1，在非对称线性算子与后台阶上可达 1.2--1.8。增广维数可以由一次训练的网络截断得到。

本文的降阶只施加于主场，辅助场仍在全网格上求解，闭合中的残差与雅可比也在全网格上组装。NE-PODG 每次查询需要计算网络候选场并构造互补空间，增广空间的维数也高于 SA-PODG，单次查询耗时的中位数在方腔自然对流上为 301 ms 对 23 ms，在 Couette--Graetz 上为 104 ms 对 5.9 ms，在后台阶上为 2.92 s 对 1.55 s。从最近训练快照出发的高保真求解在 GPU 上耗时 58 ms--5.4 s，在后台阶以外的五个算例上快于本文实现的 NE-PODG。对残差与雅可比做超降阶是降低当前全网格组装成本的一项直接途径。
